\section{Discussion}
\label{sec:discuss}
In this section, we firstly discuss two observations in preventing reward hacking and LLM's degenerate behaviors in complex end-to-end GPU programming. Then we also discuss current limitation and possible future directions.

\subsection{Format check fails to prevent reward hacking}
In our early experiment, we try to use format check to prevent reward hacking in RL, which is also applied by Kevin-32B~\cite{kevin}. However, format check fails to prevent reward hacking, both in our experiments and Kevin-32B's result in Table~\ref{tab:hack}. We attribute the failure to two reasons: First, some reward hacking behaviors are very subtle and hard to be detected by a regex-based scripts, such as hard coding. Second, format check faces a trade-off, where strict checks could classify correct solution as “hacking”, and loose checks could omit some hacking solutions. Therfore, format check alone is not a good remedy for avoiding reward hacking on KernelBench. 



\subsection{Degenerate behaviors on KernelBench}
We also observe that current LLMs clearly exhibit degenerate behaviors in complex end-to-end GPU programming tasks. For example, in many Level 2/3 tasks with multiple kernels, CUDAForge and Kevin-32B show a strong tendency to only modify “easy parts,” with only minor impact on performance. For example, only write a ReLU kernel and leave all others (e.g., conv2d and GEMM) unchanged.  And they will even replace only one of multiple ReLU kernels in the reference code with a customized kernel. In this case, the program can be correct with a very minor different performance (typically less than 5\% in our experiments). A case of this kind of code generated by Kevin is shown in Appendix~\ref{app:fear}. We attribute the degenerate behaviors to two reasons:

\textbf{RL reward design pushes the model.} According to the rule-based reward design of Kevin-32B in Appendix~\ref{kevinreward}. A correct but low-speed kernel will get higher rewards than a kernel with aggressive optimizition but has a minor error, hindering model try those advanced optimization techniques. See Appendix~\ref{app:rubric} for how StitchCUDA address it.

\textbf{Current LLMs have no confidence.} We manually track the chain of thought of these models in thinking mode to figure out the reason of degenerate behaviors. We observe a lot of similar content in their chain of thought, such as “current conv2d in PyTorch is backend with cuBLAS lib, which is hard to surpass, focus on ReLU and Linear kernels first.”. This reveals that current LLMs have no confidence to challenge those real critical parts in GPU programming. And as a result, they also lose the opportunity to utilize low-level deep optimization techniques in practice, such as fine-grained tiling and tensor cores. Based on our observations, Kevin32B never use these advanced techniques in experiments, which partially explains why its speedup decreases significantly on KernelBench level 3. 

\subsection{Limitation and future directions}
Although \textsc{StitchCUDA} demonstrates strong performance on end-to-end GPU programming tasks, several important limitations remain for this area.

\textbf{Limited generalization across hardware generations.}
Our experiments cover H200 and RTX PRO 6000, but GPU optimization is highly architecture-dependent. Code that performs well on one platform may not transfer to another due to differences in hardware specs such as shared memory capacity and tensor core support.  Therefore, it remains unclear whether an LLM-generated GPU program can consistently preserve correctness and performance across hardware generations.

\textbf{Adapting LLMs to rapidly evolving GPU knowledge.}
\textsc{StitchCUDA} uses RAG to provide agents with up-to-date CUDA/GPU documentation, which partially mitigates the knowledge staleness of pretrained LLMs. However, large-scale LLM pretraining  are difficult to keep synchronized with the fast iteration of GPU architectures stacks. New hardware features, programming interfaces, and optimization practices may appear after a model is trained, and simply retrieving documentation does not guarantee that the model can correctly interpret such knowledge and apply it to CUDA programming decisions. How to make LLM agents effectively retrieve and reason over the latest hardware knowledge remains an important future direction.

\textbf{Limitations of current benchmarks.}
Current CUDA-generation benchmarks~\cite{kernelbench,zhu2026cudabench,li2025tritonbench} mainly evaluate correctness and speedup under a small set of fixed environments. They do not systematically measure cross-platform portability, robustness to software-stack changes, deployment safety, or long-term regression behavior. Deployment-quality GPU software should also be evaluated under real application semantics, such as training stability for ML training workloads, rather than only output correctness on isolated benchmark inputs. Future benchmarks should therefore move beyond single-environment performance evaluation and incorporate portability, robustness, safety, and regression testing. At the time of this paper's publication, several concurrent and recent benchmark efforts have started to move in this direction, highlighting the growing need for more robust evaluation of LLM-generated GPU programs~\cite{computeeval,cudahercules,cudabeaver}.



\section{Conclusion}
In this paper, we propose StitchCUDA, a multi-agent framework that integrates Rubric-based reinforcement learning with the Coder agent. We identify three main challenges in LLM-based automated end-to-end GPU program generation. We improve the multi-agent framework through well-designed multi-agent decomposition and fundamentally enhance the Coder's CUDA coding ability through rubric-based agentic reinforcement learning. Comprehensive experiment results show that StitchCUDA delivers excellent performance on end-to-end GPU programming tasks, significantly outperforming simple multi-agent approaches and RL models.